\section{Training}\label{sec:training}
MorphMimic is trained directly across all bodies on refined HUMOS references,
with no neutral pretraining stage, using PPO~\cite{ppo}.

\paragraph{Optimization.}
References are sampled only from variants matching the environment's asset.
Reference-state initialization starts at frame zero with probability 0.2 and
otherwise at a random valid reference time. Training episodes last at most
1,000 control steps and terminate on falls (a 0.15\,m height setting, with
permitted foot-contact bodies), while tracking-error termination is disabled, so
this rule is distinct from the evaluation failure threshold of
Section~\ref{sec:metrics}. PPO uses 32-step rollouts, discount 0.99, GAE
parameter 0.95, policy-ratio clipping 0.2, normalized advantages, and separate
Adam optimizers with learning rates $2\times10^{-5}$ for the actor and
$10^{-4}$ for the critic. Other settings and the actual run budgets are in
Table~\ref{tab:training_parameters} (Appendix~\ref{app:reproducibility}).

\paragraph{Streamed motion sampling.}
The full-scale experiment streams per-clip files through
\texttt{GlobalClipPool}. The training manifest is deterministically partitioned
across GPU ranks, and each rank keeps a bounded resident pool containing all 128
variants of its selected clips, with environment-level motion sampling
respecting the assigned morphology. The pool is rebuilt every 256 epochs: most
slots are selected by priority, combining estimated clip difficulty and an
exploration bonus, and a fraction 0.2 is reserved for uniform random selection
to maintain rehearsal outside the highest-priority set. This is a training
curriculum, not uniform sampling of the complete corpus. Validation membership
is fixed independently of pool rebuilding and does not update sampling weights,
so resident-pool and validation metrics measure different populations. Sampling
constants and the pool size actually used are in
Appendices~\ref{app:reproducibility} and~\ref{app:history}.
